% !TeX root = ../tesis.tex

% =============================================================================
%  ANEXO D — Implementaciones Computacionales de Indicadores
%  Prototipo teórico — sustituir con código de producción verificado
% =============================================================================

Este anexo reúne los prototipos de implementación en Python de cada indicador
definido en §\ref{sec:indicadores-formulas}. El código aquí presentado
corresponde a implementaciones teóricas de referencia; la versión de producción
---validada sobre los datos reales de la \zmvm{}--- se desarrollará en
investigaciones posteriores o en la fase de cierre del presente trabajo.

% -----------------------------------------------------------------------
\section{Cobertura de la Red \texorpdfstring{$C$}{C}}
\label{anx:impl-cobertura}
% -----------------------------------------------------------------------

Implementación de referencia del indicador $C$ definido en §
\ref{subsec:indicador-cobertura}, utilizando \textbf{GeoPandas} y
\textbf{Shapely} sobre los datos \gtfs{} de la red metropolitana.

\begin{lstlisting}[style=estiloPython,
  caption={Cálculo del indicador $C$ con GeoPandas (prototipo).},
  label={lst:cobertura}]
# GeoPandas + Shapely
gdf_stops = gpd.read_file("paradas_red.geojson")
buffers = gdf_stops.geometry.buffer(400)       # Buffer 400 m
covered_area = buffers.unary_union             # Union de poligonos
city_polygon = gpd.read_file("zmvm.geojson").geometry[0]
C = covered_area.intersection(city_polygon).area / city_polygon.area
\end{lstlisting}

% -----------------------------------------------------------------------
\section{Alimentación Capilar --- Fuerza Nodal \texorpdfstring{($C_i$)}{(Ci)}}
\label{anx:impl-capilar}
% -----------------------------------------------------------------------

Implementación de referencia del indicador $C_i$ definido en §
\ref{subsec:alimentacion-capilar}, utilizando \textbf{NetworkX}. El código
calcula la fuerza nodal como suma de pesos de aristas de superficie, conforme a
la fórmula $C_i = \sum_{j \in S} A_{ij} \cdot w_{ij}$.

\begin{lstlisting}[style=estiloPython,
  caption={Cálculo del indicador $C_i$ con NetworkX (prototipo).},
  label={lst:fuerza-nodal}]
# NetworkX - Fuerza nodal C_i: suma de pesos de aristas de superficie
capilar_strength = {
    node: sum(d.get('weight', 0)
              for _, _, d in G.edges(node, data=True)
              if d.get('tipo') == 'superficie')
    for node in G.nodes()
}
\end{lstlisting}

% -----------------------------------------------------------------------
\section{\textit{Detour Factor} --- Índice de Ruta Directa \texorpdfstring{($DI$)}{(DI)}}
\label{anx:impl-detour}
% -----------------------------------------------------------------------

Implementación de referencia del indicador $DI$ definido en §
\ref{subsec:ruta-directa}. Usa \textbf{geopy} para la distancia haversine y
\textbf{NetworkX} para el camino mínimo ponderado por distancia en km.

\begin{lstlisting}[style=estiloPython,
  caption={Cálculo del \textit{Detour Factor} $DI$ (prototipo).},
  label={lst:detour}]
from geopy.distance import geodesic
import networkx as nx

d_euclidiana = geodesic(coords_origen, coords_destino).km
ruta_minima = nx.shortest_path(G, origen, destino, weight='distancia_km')
d_red = sum(G[u][v]['distancia_km'] for u, v in zip(ruta_minima, ruta_minima[1:]))
DI = d_red / d_euclidiana
\end{lstlisting}

% -----------------------------------------------------------------------
\section{Centralidad de Intermediación (\textit{Betweenness Centrality}) \texorpdfstring{($B(v)$)}{(B(v))}}
\label{anx:impl-betweenness}
% -----------------------------------------------------------------------

Implementación de referencia del indicador $B(v)$ definido en §
\ref{subsec:centralidad-intermediacion}, utilizando \textbf{NetworkX} con peso
por tiempo de viaje real y normalización en $[0, 1]$.

\begin{lstlisting}[style=estiloPython,
  caption={Cálculo de \textit{Betweenness Centrality} $B(v)$ con NetworkX (prototipo).},
  label={lst:betweenness}]
import networkx as nx
centrality = nx.betweenness_centrality(G, weight='tiempo_minutos', normalized=True)
# Ordenar por centralidad descendente
top_nodes = sorted(centrality.items(), key=lambda x: x[1], reverse=True)[:10]
\end{lstlisting}

% -----------------------------------------------------------------------
\section{Tiempo de Viaje Promedio \texorpdfstring{($T$)}{(T)}}
\label{anx:impl-tiempo}
% -----------------------------------------------------------------------

Implementación de referencia del indicador $T$ definido en §
\ref{subsec:tiempo-viaje}, utilizando \textbf{NetworkX}. Para redes de gran
escala se recomienda sustituir el cómputo exhaustivo por una muestra
representativa de pares O--D.

\begin{lstlisting}[style=estiloPython,
  caption={Cálculo del Tiempo de Viaje Promedio $T$ con NetworkX (prototipo).},
  label={lst:tiempo}]
import networkx as nx
# Para redes grandes: sustituir por muestra representativa de pares O-D
T = nx.average_shortest_path_length(G, weight='tiempo_minutos')
\end{lstlisting}

% -----------------------------------------------------------------------
\section{Robustez Geométrica \texorpdfstring{($\Delta E$)}{(Delta E)}}
\label{anx:impl-robustez}
% -----------------------------------------------------------------------

Implementación de referencia del indicador $\Delta E$ definido en §
\ref{subsec:robustez-geometrica}. Pendiente de validación empírica sobre la red
real de la \zmvm{} (véase nota aclaratoria en §\ref{sec:indicadores-formulas}).

\begin{lstlisting}[style=estiloPython,
  caption={Cálculo de Robustez Geométrica $\Delta E$ con NetworkX (prototipo).},
  label={lst:robustez}]
import networkx as nx, copy

def eficiencia_global(G):
    T = nx.average_shortest_path_length(G, weight='tiempo_minutos')
    return 1 / T

E0 = eficiencia_global(G)

# Remover nodo critico (reemplazar con nodo real de mayor B(v))
G_fallido = copy.deepcopy(G)
G_fallido.remove_node('Pantitlan')
Ef = eficiencia_global(G_fallido)

delta_E = (E0 - Ef) / E0
print(f"Caida de eficiencia: {delta_E:.2%}")
\end{lstlisting}

% -----------------------------------------------------------------------
\section{Penalización por Transferencia \texorpdfstring{($W_{s,h}$)}{(W)}}
\label{anx:impl-transferencia}
% -----------------------------------------------------------------------

Implementación de referencia del indicador $W_{s,h}$ definido en §
\ref{subsec:penalizacion-transferencia}. El costo de transbordo suma el tiempo
de caminata entre estaciones más la espera promedio hasta el siguiente servicio.
La implementación de producción está integrada en \texttt{graph\_builder.py}; el
prototipo aquí presentado es de referencia.

\begin{lstlisting}[style=estiloPython,
  caption={Cálculo de la Penalización por Transferencia $W_{s,h}$ (prototipo).},
  label={lst:transferencia}]
# Pendiente: integracion completa en graph_builder.py
# Costo de transbordo = tiempo_caminata + espera_promedio
# espera_promedio = frecuencia_minutos / 2

def costo_transbordo(distancia_m, velocidad_peatonal_kmh, frecuencia_min):
    tiempo_caminata = (distancia_m / 1000) / velocidad_peatonal_kmh * 60
    espera = frecuencia_min / 2
    return tiempo_caminata + espera
\end{lstlisting}
